\documentclass[10pt,a4paper]{amsart}
\usepackage[margin=19mm]{geometry}
\usepackage{amsmath,amssymb,mathtools}
\usepackage[scr=rsfs,cal=euler]{mathalfa}
\usepackage{xfrac}
\usepackage[unicode,hidelinks,bookmarks=false]{hyperref}
\newcommand{\ie}{i.e.\ }
\newcommand{\cf}{cf.\ }
\newcommand{\resp}{respectively\ }
\newcommand{\Subsection}{\subsection}
\providecommand{\nobreakdash}{\nobreak\hbox{-}}
\setlength{\emergencystretch}{2em}
\multlinegap0pt
\usepackage{bm}
\usepackage{bbm}
\usepackage{dsfont}
\usepackage{xspace}
\usepackage{cancel}
\usepackage{mathtools}
\usepackage{amsthm}
\makeatletter
\@ifundefined{theorem}{\newtheorem{theorem}{Theorem}}{}
\@ifundefined{proposition}{\newtheorem{proposition}[theorem]{Proposition}}{}
\makeatother
\newcommand{\vect}[1]{\boldsymbol{#1}}
\title[Covariant Lyapunov vectors as global solutions of a partial ]{Covariant Lyapunov vectors as global solutions of a partial differential equation on the phase space}
\author{Massimo Marino, Doriano Brogioli}
\date{Source: arXiv:2301.12543v3 (CC BY 4.0)}
\begin{document}
\maketitle

\noindent\emph{Source and licence.} Covariant Lyapunov vectors as global solutions of a partial differential equation on the phase space. Version arXiv:2301.12543v3 (CC BY 4.0), distributed under the Creative Commons Attribution 4.0 International licence, as recorded in the arXiv licence metadata for this exact version and shown on the abstract page https://arxiv.org/abs/2301.12543v3. Full rights notice: licenses/arxiv-2301.12543-CC-BY-4.0.txt.\par\smallskip

\noindent\textit{Editorial scope.} Counted unit: the Proposition whose source label is \texttt{prop:geo:foliation:CLF} in the article (source0.tex lines 1443--1459, source label \texttt{prop:geo:foliation:CLF}), delivered together with the article's own complete original proof of it (source0.tex lines 1461--1544, one \texttt{proof} environment of 84 lines). No mathematics is written, repaired or re-derived: statement and proof are the article's own text line for line. Required context retained uncounted: eq:xp (lines 197--201); eq:alt:def (lines 700--703); eq:walt:def (lines 764--766); prop:1 (lines 831--855); prop:alt:def (lines 900--922); prop:blim (lines 1173--1188). Every definition, display and label of those lines is retained unchanged. Omitted from this package and not redistributed: 1--196, 202--699, 704--763, 767--830, 856--899, 923--1172, 1189--1442, 1460--1460, 1545--1773. Nothing inside a retained excerpt is omitted or altered: every retained body is delivered exactly as the article's lines are written, so no omission receipt of any kind accompanies this package. Citations: the retained excerpts carry no citation command at all, so no bibliography entry of the article is transcribed and no reference is redistributed. Reference adaptation: none was needed and none was made; every reference inside the delivered unit resolves inside it, so no unresolved reference or citation is produced. Printed slips kept literal: no printed word, symbol, hypothesis or number of the counted statement or of its proof was altered. Layout and class: the article's own class is \texttt{article} with packages including graphicx, amsfonts, amsmath, amssymb, amsthm, natbib, xcolor; the wrapper is the lane's pinned \texttt{amsart} editorial wrapper at 10pt on A4, which restates the theorem-like environments the retained text needs and adds the article's own macro definitions that the retained text needs (\texttt{\textbackslash vect}), with extra wrapper packages (bm, bbm, dsfont, xspace, cancel, mathtools). The delivered numbering is the unit's own. Assets: no retained span contains a figure environment, a graphics-inclusion command, a PDF-page inclusion, a table or a TikZ picture; no image file is copied into this package, the package declares no asset dependency and no image file is redistributed with it. Licence: version arXiv:2301.12543v3 of this article is distributed under the Creative Commons Attribution 4.0 International licence, as recorded in the arXiv licence metadata for this exact version and shown on the abstract page https://arxiv.org/abs/2301.12543v3. A full rights notice is delivered at licenses/arxiv-2301.12543-CC-BY-4.0.txt. Characters: every retained excerpt is pure ASCII, so the delivered engine, XeLaTeX with its default Latin Modern text font, reports no missing character.
\begin{equation}
  \frac{\mathrm{d}}{\mathrm{d}t} x^{\mu}(t) = 
  F^{\mu}\left[\vect{x}(t)\right] 
  \label{eq:xp}
\end{equation}

 \begin{equation}
    \mathcal{L}_{\vect{F}} \vect{v} + b \vect v = 0 \,.
    \label{eq:alt:def}
  \end{equation}  

\begin{equation}\label{eq:walt:def}
\mathcal{L}_{\vect{F}}\vect{w} + c \vect w = 0
\end{equation}

\begin{proposition}  \label{prop:1}
Let $\mu$ be a positive measure on the Riemaniann manifold $X$ such
that $\mu(X)<+\infty$, and let the vector field $\vect{F}$ generate an
ergodic flow on $X$ which preserves the measure $\mu$.  Let $\vect{v}$
be a CLF, corresponding to a nondegenerate LE $\lambda$, on an
invariant domain \mbox{$D\subseteq X$} such that $\mu(D)=\mu(X)$.
Then the Lie derivative of $\vect{v}$ along $\vect{F}$ exists, and
there exists a scalar field $b$ such that the differential equation 
\[
 \mathcal{L}_{\vect{F}} \vect{v} + b \vect v = 0 
\]
holds on $D$. 
Let us also suppose that the function 
$\ln \lVert \vect{v}\left(\vect {x}\right) \rVert$ is integrable over $X$ 
with respect to the measure $\mu$. Then
  \begin{equation}\label{eq:lambda_b}
    \lambda=\left<b\right>\,, 
  \end{equation}
  where $\left<b\right>$ denotes the
  average of $b$ over the manifold $X$:
  \begin{equation}\label{eq:aver_b}
    \left<b\right>=  \frac {1}{\mu(X)}
    \int_X \mathrm{d}\mu (\vect{x})\, b(\vect{x}) \, .
  \end{equation}
\end{proposition}

\begin{proposition}
  \label{prop:alt:def}
Let $\mu$ be a positive measure on the Riemaniann manifold $X$ such
that $\mu(X)<+\infty$, and let the vector field $\vect{F}$ generate an
ergodic flow on $X$ which preserves the measure $\mu$.  Let $\vect{v}$
and $b$ be respectively a nonvanishing vector field and a scalar
function satisfying the equation
\[
 \mathcal{L}_{\vect{F}} \vect{v} + b \vect v = 0 
\]
on an invariant domain \mbox{$D\subseteq X$}, such that
$\mu(D)=\mu(X)$. Let us suppose that the function
\begin{equation}\label{eq:bb2}
c(\vect{x})=b(\vect{x})+\mathcal{L}_{\vect{F}}
\ln \lVert \vect{v}\left(\vect {x}\right) \rVert 
\end{equation}
is integrable over $X$ with respect to the measure $\mu$, and
that $\lVert \vect{v}\left(\vect {x}\right) \rVert$ is differentiable
on $D$. Then $\vect{v}$ is a CLF with LE $\lambda=\left<c\right>$ on
an invariant domain $D'\subseteq D$ such that $\mu(D')=\mu(X)$. If, in
addition, also $\ln \lVert \vect{v}\left(\vect {x}\right) \rVert$ is
integrable over $X$, then $\left<b\right>=\left<c\right>=\lambda$.
\end{proposition}

\begin{proposition}\label{prop:blim}
Let $\vect{F}$ be a differentiable vector field on the Riemaniann manifold
$X$, and let $D\subseteq X$ be an invariant set with respect to the evolution
generated by $\vect{F}$. Let $\vect{w}$
be a CLF, corresponding to a nondegenerate LE $\lambda$, on an
invariant domain \mbox{$D\subseteq X$} such that $\mu(D)=\mu(X)$. If
\begin{equation}\label{eq:wnorm1}
\lVert \vect{w}(\vect{x})\rVert=1 \ \forall\,\vect{x}\in D\,,
\end{equation}
then $\vect w$
satisfies Eq.~(\ref{eq:walt:def}) on the domain $D$ with 
\begin{equation}\label{eq:blie}
c= \frac 12
w^\mu \left(\mathcal{L}_{\vect{F}} g_{\mu\nu}\right) w^\nu\,.
\end{equation}
\end{proposition}

\begin{proposition} \label{prop:geo:foliation:CLF}
Let $\mu$ be a positive measure on the Riemaniann manifold $X$
such that $\mu(X)<+\infty$, and
let the vector field $\vect{F}$ generate, according to Eq.~(\ref{eq:xp}), 
an ergodic flow on $X$ which preserves the measure $\mu$.
Let $\vect{F}$ be a CLF with LE $\lambda=0$ and be such that
the function $\lVert\mathcal{L}_{\vect{F}} g(\vect{x})\rVert$
is integrable over $X$ with respect to the measure $\mu$.
Let also
  $\vect{v}_i(\vect{x})$, for $i=1, \dots, n-1$, 
be $n-1$ additional CLFs, on an invariant domain
  $D\subseteq X$, corresponding to nondegenerate LEs $\lambda_i$.
  If a 2-dimensional foliation $\Phi$ of $D$ is such that each leave contains whole
  trajectories of $\vect{F}$, then there exists one index $\bar i$,
with $1\leq \bar i\leq n-1$, such that the foliation is generated by 
the couple $(\vect{F},\vect{v}_{\bar i})$.
\end{proposition}

\begin{proof}
At each point $\vect{x}\in D$ of a 2-dimensional leave 
of $\Phi$ the vector $\vect{F}(\vect{x})$ is tangent to the leave, 
since the leave contains whole trajectories. We can then take,
as a  basis of the tangent space of the leave, the vector
$\vect{F}(\vect{x})$ and a vector $\vect{W}(\vect{x})$ which can be
  expressed as a linear combination of the $n-1$ CLVs 
$\vect{v}_1(\vect{x})$, $\dots$,  $\vect{v}_{n-1}(\vect{x})$.
If we also impose the condition
$\lVert \vect{W}(\vect{x})\rVert =1$, then the vector $\vect{W}
(\vect{x})$ is univocally determined (up to the sign)
at each $\vect{x}\in X$, and we
obtain in this way a vector field on $D$:
  \begin{equation} \label{eq:geo:def:W}
    \vect{W}(\vect{x}) = \sum_{i=1}^{n-1} c_i (\vect{x})
\vect{v}_i (\vect{x})\,.
  \end{equation}
  
Since $\Phi$ is a 2-dimensional foliation, 
according to Frobenius theorem $\vect{F}$ and $\vect{W}$ must be
  two involutive vector fields, thus the Lie derivative
  $\mathcal{L}_{\vect{F}}\vect{W}$ must be a linear combination of
  $\vect{F}$ and $\vect{W}$. This means that
  \begin{equation}\label{eq:lie_fw}
    \mathcal{L}_{\vect{F}}\vect{W} = \alpha \vect{F} - \beta \vect{W}\,,
  \end{equation}
  where $\alpha$ and $\beta$ are two scalar fields, and
using Eq.~(\ref{eq:geo:def:W}) to express $\vect{W}$ we get
  \begin{equation}\label{eq:prop_geom}
    \sum_{i=1}^{n-1} \mathcal{L}_{\vect{F}}(c_i \vect{v}_i) =
    \alpha \vect{F} + \beta \sum_{i=1}^{n-1} c_i \vect{v}_i\,.
  \end{equation}
It follows from proposition \ref{prop:1} that for any $i=1, \dots, n-1$
there exists a scalar function $b_i$ such that the equation
 \begin{equation}
    \mathcal{L}_{\vect{F}} \vect{v}_i + b_i \vect v_i = 0 
    \label{eq:alt:defi}
  \end{equation}  
holds on $D$. From Eq.~(\ref{eq:prop_geom}) we then get
  \begin{equation}\label{eq:foliation_c}
\alpha \vect{F} -\sum_{i=1}^{n-1}\left[\mathcal{L}_{\vect{F}}c_i-
(b_i-\beta)c_i \right] \vect{v}_i = 0    \,.
  \end{equation}
Since the set of vectors $\left(\vect{F},\vect{v}_1, \dots, 
\vect{v}_{n-1}\right)$ is
linearly independent at all points $\vect{x}\in D$, 
from the above equation we get that for all $\vect{x} \in D$
\begin{align}
\alpha(\vect{x})&=0 \,, \label{eq:alpha}\\
 \mathcal{L}_{\vect{F}}c_i(\vect{x})&=\left[b_i(\vect{x})-
\beta(\vect{x})\right]c_i(\vect{x})\ \forall\,i=1,\dots, n-1\,. \label{eq:coeffci}
\end{align}
According to Eq.~(\ref{eq:alpha}) we can rewrite Eq.~(\ref{eq:lie_fw}) as
 \begin{equation}\label{eq:lie_fw1}
    \mathcal{L}_{\vect{F}}\vect{W} +\beta \vect{W}=0\,,
  \end{equation}
which has the same form as Eq.~(\ref{eq:alt:def}) with $b=\beta$.
Since $\lVert \vect{W}(\vect{x})\rVert =1$, by applying proposition 
\ref{prop:blim} we obtain
\begin{equation}%\label{eq:blim}
\lvert \beta(\vect{x})\rvert \leq \frac 12 \lVert 
\mathcal{L}_{\vect{F}} g(\vect{x}) \rVert
\ \forall\, \vect{x}\in X\,.
\end{equation}
Since by hypothesis the function on the right-hand side
is integrable over $X$, the above equation
implies that also $\beta$ is integrable. We can then apply proposition
\ref{prop:alt:def} and deduce from Eq.~(\ref{eq:lie_fw1}) that
$\vect{W}$ is a CLF with LE $\lambda=\left<\beta\right>$. But since 
$\left(\vect{v}_1, \dots, \vect{v}_{n-1}\right)$ is a
set of nondegenerate CLFs, it follows from Eq.~(\ref{eq:geo:def:W})
that there must be only one index $\bar i$,
with $1\leq \bar i\leq n-1$, such that the function $c_{\bar i} (\vect{x})$
is not identically 0. Hence 
\begin{equation} \label{eq:foliation_w}
    \vect{W}(\vect{x}) = c_{\bar i} (\vect{x})
\vect{v}_{\bar i} (\vect{x})
  \end{equation}
and $\left<\beta\right>=\lambda_{\bar i}$. 

Since by construction the couple $(\vect{F},\vect{W})$ generates the
foliation $\Phi$, it follows from Eq.~(\ref{eq:foliation_w}) that
also the couple $(\vect{F},\vect{v}_{\bar i})$ generates $\Phi$.
\end{proof}

\end{document}
